\scanpage{59}

% manuscript line(s) 1–3
\deleted{Trois} \inserted{Quatre} catégories fibrées sur $(\mathrm{Sch}_{\unsure{?}})$ (plus gén. sur la catégorie des topos loc. annelés $\ldots$), canoniquement équivalentes
% manuscript line(s) 4–11
\diagram{\begin{tikzcd}[ampersand replacement=\&,column sep=small,row sep=small] \text{\textcircled{1}}\arrow[d,leftrightarrow]\&\parbox{.78\linewidth}{Catégorie des rev. quadratiques $X'/X$ de $X\in\mathrm{Ob}\,\mathrm{Sch}$,}\\ \text{\textcircled{$\deleted{\unsure{2}}1'$}}\arrow[d]\&\parbox{.78\linewidth}{Catégorie des alg \inserted{$\mathcal A$} commutatives loc libres de rang $2$ sur $\underline O_X$}\\ \text{\textcircled{2}}\&\parbox{.78\linewidth}{Catégorie des triples $(\mathcal A,N,e)$, où $(\mathcal A,N)$ est un module quadratique loc. libre de rang $2$ sur $\underline O_X$, et où $e$ est une section de $\mathcal A$ telle que \(\displaystyle N(e)=1\)}\\ \text{\textcircled{3}}\arrow[uuu,bend left=35]\&\parbox{.78\linewidth}{Catégorie des triples $(L,P,Q)$, où $L$ est un $\underline O_X$-Module inv. \inserted{(identifié au schéma en droites homogène corresp)}, $P$ un $L$-torseur, et $Q$ une application polynomiale}\end{tikzcd}}
% manuscript line(s) 12
\[
P\longrightarrow L^{\otimes2}
\]
% manuscript line(s) 13–14
dont la partie homogène de degré $2$ soit $u\unsure{\mathbin{-}|\mapsto}u^2$, i.e. on suppose que
% manuscript line(s) 15
\[
Q(x+u)=Q(x)+u^2+\mathrm{Lin}_u(x)
\]
% manuscript line(s) 16–18
où $\mathrm{Lin}_u:P_{S}\mathbin{-}L_S^{\otimes2}$ est une \unsure{application|appl.} \deleted{\unsure{linéaire|lin.}} \uline{affine} (quand $u$ est une section loc.\textsuperscript{le} \inserted{fppf, disons} sur $S/X$)

% manuscript line(s) 19–22
$\updownarrow$\quad\textcircled{$3'$} Catégorie des triples $(0\to L\overset{i}\to E\overset{\psi}\to\underline O_X\to0,Q)$ d'une suite exacte $0\mathbin{-}L\mathbin{-}E\mathbin{-}\underline O_X\mathbin{-}0$, avec $L$ inv., et d'une forme quadratique $Q$ sur $E$ à val. dans $L^{\otimes2}$, dont la restriction à $L$ soit $u\mathbin{-}u^{\otimes2}$.

% manuscript line(s) 23–25
On va transformer \textcircled{$3'$}, en \deleted{\unsure{pl}} considérant $L^{\otimes2}$ comme plongé dans $\mathrm{Sym}^2(E)$ ($\simeq\mathrm{Quad}(\mathcal A)$, où $\mathcal A\overset{\mathrm{def}}=\check E$), lequel a une suite de composition
% archive sheet number: 58
